% !TEX root = TRENCH_DIFFEQ.tex

\usepackage{mathtools}
\usepackage{lua-unicode-math}
\usepackage[usenames, dvipsnames, svgnames, table]{xcolor}
%\usepackage[T1]{fontenc}
%\usepackage{textcomp}
%\renewcommand{\rmdefault}{ptm}
%\usepackage[scaled=0.92]{helvet}
%\usepackage[psamsfonts]{amsfonts}
\usepackage{amsmath, amsbsy}
\usepackage{calc}
\usepackage{booktabs}
\usepackage[per-mode = symbol]{siunitx}
\usepackage{capt-of}
\usepackage{zref-clever}
\usepackage{etoolbox}
\usepackage{pgfplots}
\pgfplotsset{compat=1.18}

\usepackage[%dvips,
 bookmarks = true,
 colorlinks=true,
 plainpages = false,
 citecolor = blue,
 urlcolor = blue,
 filecolor = blue,
 pageanchor = true]
{hyperref}

\zcRefTypeSetup{chapter}{cap}
\zcRefTypeSetup{definition}{cap}
\zcRefTypeSetup{equation}{cap}
\zcRefTypeSetup{example}{cap}
\zcRefTypeSetup{figure}{cap}
\zcRefTypeSetup{section}{cap}
\zcRefTypeSetup{table}{cap}
\zcRefTypeSetup{theorem}{cap}
\zcRefTypeSetup{enumXi}{
  Name-sg=Exercise,
  Name-pl=Exercises,
  name-sg=exercise,
  name-pl=exercises,
  rangesep={\textendash},
  cap
}
\zcRefTypeSetup{enumXvii}{
  Name-sg=Exercise,
  Name-pl=Exercises,
  name-sg=exercise,
  name-pl=exercises,
  rangesep={\textendash},
  cap
}

\newtoggle{boundaryvalues}

\EditInstance{heading}{section}{title-decls=\raggedright}

\DeclareMathOperator{\tr}{tr}

\DeclareSIUnit\dyne{\text{dyn}}
\DeclareSIUnit[quantity-product = {}]\degreeFahrenheit{\text{\textdegree F}}
\DeclareSIUnit\feet{\text{ft}}
\DeclareSIUnit\foot{\text{ft}}
\DeclareSIUnit\gallon{\text{gal}}
\DeclareSIUnit\pound{\text{lb}}

\newcommand{\vect}[1]{\mathbf{#1}}
\newcommand{\ivect}{\vect{i}}
\newcommand{\jvect}{\vect{j}}

\newcommand{\laplace}{{\mathcal L}}

\DeclareTextCommandDefault{\textRightarrow}{\ensuremath{\Rightarrow}}

\newcommand{\amatrix}[2]{
 \left[
  \begin{matrix}#1\end{matrix}
 \hspace{\tabcolsep}\middle|\hspace{\tabcolsep}
  \begin{matrix}#2\end{matrix}
 \right]
}

\ExplSyntaxOn
% TODO?
\keys_define:nn { zref-clever/reference }
  {
    CAP .bool_set:N = \l_tmpa_bool,
    CAP .default:n = false
  }
\ExplSyntaxOff

\usepackage{xstring}

\newcommand*{\personHref}[2]{%
 \IfStrEqCase{#1}{%
  {Libby}{\href{https://www.nobelprize.org/prizes/chemistry/1960/libby/facts/}{\color{blue}\slshape#2}}%
  {Malthus}{\href{http://en.wikipedia.org/wiki/Thomas_Robert_Malthus}{\color{blue}\slshape#2}}%
 }[\href{https://mathshistory.st-andrews.ac.uk/Biographies/#1/}{\color{blue}\slshape#2}]%
}



%\colorlet{blue}{DarkBlue}
\colorlet{red}{DarkRed}

\setcounter{secnumdepth}{2}
\usepackage[a4paper]{geometry}
\usepackage{float}
\usepackage{upref}

\usepackage{array}
\newcolumntype{H}{>{\setbox0=\hbox\bgroup}c<{\egroup}@{}}

\usepackage{graphicx}
\graphicspath{%
    {converted_graphics/}% inserted by PCTeX
    {EPS/}% inserted by PCTeX
}
%\usepackage{epstopdf}
%\epstopdfsetup{update}

\numberwithin{example}{section}
\numberwithin{equation}{section}
\numberwithin{theorem}{section}
\numberwithin{table}{section}
\numberwithin{figure}{section}

%\newcommand{\place}{\bigskip\hrule\bigskip\noindent}
%\newcommand{\threecol}[3]{\left[\begin{array}{r}#1\\#2\\#3\end{array}\right]}
\newcommand{\threecol}[3]{\begin{bmatrix}#1\\#2\\#3\end{bmatrix}}
%\newcommand{\threecolj}[3]{\left[\begin{array}{r}#1\\[1\jot]#2\\[1\jot]#3\end{array}\right]}
\newcommand{\threecolj}{\threecol}
\newcommand{\lims}[2]{\,\bigg|_{#1}^{#2}}
%\newcommand{\twocol}[2]{\left[\begin{array}{l}#1\\#2\end{array}\right]}
\newcommand{\twocol}[2]{\begin{bmatrix}#1\\#2\end{bmatrix}}
%\newcommand{\ctwocol}[2]{\left[\begin{array}{c}#1\\#2\end{array}\right]}
\newcommand{\ctwocol}{\twocol}
%\newcommand{\cthreecol}[3]{\left[\begin{array}{c}#1\\#2\\#3\end{array}\right]}
\newcommand{\cthreecol}{\threecol}
\newcommand{\eqline}[1]{\centerline{\hfill$\displaystyle#1$\hfill}}
%\newcommand{\twochar}[4]{\left|\begin{array}{cc}
%#1-\lambda&#2\\#3&#4-\lambda\end{array}\right|}
\newcommand{\twochar}[4]{\begin{vmatrix}#1-\lambda&#2\\#3&#4-\lambda\end{vmatrix}}
%\newcommand{\twobytwo}[4]{\left[\begin{array}{rr}
%#1&#2\\#3&#4\end{array}\right]}
\newcommand{\twobytwo}[4]{\begin{bmatrix}#1&#2\\#3&#4\end{bmatrix}}
%\newcommand{\threechar}[9]{\left[\begin{array}{ccc}
%#1-\lambda&#2&#3\\#4&#5-\lambda&#6\\#7&#8&#9-\lambda\end{array}\right]}
\newcommand{\threechar}[9]{\begin{bmatrix}#1-\lambda&#2&#3\\#4&#5-\lambda&#6\\#7&#8&#9-\lambda\end{bmatrix}}
%\newcommand{\threebythree}[9]{\left[\begin{array}{rrr}
%#1&#2&#3\\#4&#5&#6\\#7&#8&#9\end{array}\right]}
\newcommand{\threebythree}[9]{\begin{bmatrix}#1&#2&#3\\#4&#5&#6\\#7&#8&#9\end{bmatrix}}
% TODO
%\newcommand{\solutionpart}[1]{\bigskip\noindent{%\underbar
%\color{blue}\scshape Solution({\bfseries #1})\ }
%}
\newcommand{\Cex}{\fbox{\textcolor{red}{C}}\, }
\newcommand{\CGex}{\fbox{\textcolor{red}{C/G}}\, }
\newcommand{\Lex}{\fbox{\textcolor{red}{L}}\, }
%\newcommand{\matfunc}[3]{\left[\begin{array}{cccc}#1_{11}(t)&#1_{12}(t)&\cdots
%&#1_{1#3}(t)\\#1_{21}(t)&#1_{22}(t)&\cdots&#1_{2#3}(t)\\\vdots&
%\vdots&\ddots&\vdots\\#1_{#21}(t)&#1_{#22}(t)&\cdots&#1_{#2#3}(t)
%\end{array}\right]}
\newcommand{\matfunc}[3]{\begin{bmatrix}#1_{11}(t)&#1_{12}(t)&\cdots
&#1_{1#3}(t)\\#1_{21}(t)&#1_{22}(t)&\cdots&#1_{2#3}(t)\\\vdots&
\vdots&\ddots&\vdots\\#1_{#21}(t)&#1_{#22}(t)&\cdots&#1_{#2#3}(t)
\end{bmatrix}}

%\newcommand{\col}[2]{\left[\begin{array}{c}#1_1\\#1_2\\\vdots\\#1_#2\end{array}\right]}
\newcommand{\col}[2]{\begin{bmatrix}#1_1\\#1_2\\\vdots\\#1_#2\end{bmatrix}}
%\newcommand{\colfunc}[2]{\left[\begin{array}{c}#1_1(t)\\#1_2(t)\\\vdots\\#1_#2(t)\end{array}\right]}
\newcommand{\colfunc}[2]{\begin{bmatrix}#1_1(t)\\#1_2(t)\\\vdots\\#1_#2(t)\end{bmatrix}}

%\newcommand{\cthreebythree}[9]{\left[\begin{array}{ccc}
%#1&#2&#3\\#4&#5&#6\\#7&#8&#9\end{array}\right]}
\newcommand{\cthreebythree}[9]{\begin{bmatrix}#1&#2&#3\\#4&#5&#6\\#7&#8&#9\end{bmatrix}}

\newcommand{\technology}{\bigskip\noindent\fbox{\color{red}USING
TECHNOLOGY}\par}

\newcommand{\bigbox}[1]{\vskip10pt\vbox{\hrule\hbox{\vrule\kern5pt
\vbox{\noindent
\strut\kern5pt\mbox{}\newline#1\mbox{}\newline
\kern5pt\strut}\kern5pt\vrule}\hrule}}

\newcommand{\blankbox}[1]{\vskip10pt\vbox{\hbox{\kern5pt
\vbox{\noindent\strut\kern5pt\vspace*{-15pt}#1\vspace*{-20pt}
\kern5pt\strut}\kern5pt}}}

%\newenvironment{newtable}[1]{ \begin{quote}
%Table\space\refstepcounter{table}\thetable. #1}{\end{quote}}

\newcommand{\sectionanswers}[1]{%
 \ParseLaTeXeHeading{section}{\BooleanTrue}{
  heading-decls = \normalfont\normalsize\bfseries,
  after-vspace = \medskipamount,
  before-vspace = \medskipamount,
  bookmark = {Section #1 Answers},
%  toc = {Section #1 Answers},
  running = {Section #1 Answers}
 }{Section #1 Answers, pp.~\pageref{exer:#1.1}\thinspace ff}%
}

\newcommand{\indexterm}[1]{\emph{\color{blue}#1}\index{#1}}

\setlength{\oddsidemargin}{15.5pt}
\setlength{\evensidemargin}{15.5pt}

\usepackage{makeidx}
\makeindex

\title{Elementary Differential Equations}
\author{William F. Trench}
